% ============================================================
%  Agenda / Bitacora de avance - TFG AutoPipe
%  Autor: Ricardo Antonio Borbon Mena
%
%  COMO USAR ESTE ARCHIVO
%  ----------------------
%  1. Para registrar un dia, agregue UNA linea al final de la
%     seccion "Bitacora", con este formato:
%
%         \entrada{AAAA-MM-DD}{horas}{que hice}{bloqueos o decisiones}
%
%     Ejemplo:
%         \entrada{2026-08-07}{5}{Implemente el cargador tabular.}{Sin bloqueos.}
%
%  2. Las horas admiten decimales: {4.5} es valido.
%  3. El acumulado de horas se calcula solo; no hay que actualizarlo a mano.
%  4. Para separar semanas, agregue: \semana{3}{2026-08-17 a 2026-08-23}
%  5. Compilar con:  pdflatex Agenda.tex   (dos veces si cambia el indice)
% ============================================================

\documentclass[11pt,letterpaper]{article}

\usepackage[utf8]{inputenc}
\usepackage[T1]{fontenc}
\usepackage[spanish,es-tabla]{babel}
\usepackage{lmodern}
\usepackage[letterpaper,top=2.2cm,bottom=2.2cm,left=2.5cm,right=2.5cm]{geometry}
\usepackage{xcolor}
\usepackage{booktabs}
\usepackage{array}
\usepackage{enumitem}
\usepackage{xfp}
\usepackage{fancyhdr}

\definecolor{tecnavy}{HTML}{003366}
\definecolor{gris}{HTML}{666666}

\setlength{\parindent}{0pt}
\setlength{\parskip}{4pt}
\renewcommand{\arraystretch}{1.2}

\pagestyle{fancy}
\fancyhf{}
\renewcommand{\headrulewidth}{0.4pt}
\fancyhead[L]{\small\textit{AutoPipe}}
\fancyhead[R]{\small Bitácora de avance}
\fancyfoot[C]{\thepage}

% --- Presupuesto total del proyecto -------------------------
\newcommand{\presupuesto}{336}

% --- Acumulador de horas (se actualiza solo) ----------------
\newcommand{\horasacum}{0}

% --- Una entrada de bitácora --------------------------------
% #1 fecha  #2 horas  #3 qué hice  #4 bloqueos/decisiones
\newcommand{\entrada}[4]{%
  \xdef\horasacum{\fpeval{\horasacum + #2}}%
  \par\medskip
  \noindent\textbf{\color{tecnavy}#1}\quad\textbf{#2\,h}%
  \hfill{\small\color{gris}acumulado: \horasacum\,h / \presupuesto\,h}\par
  \nopagebreak
  \noindent\textit{Hice:} #3\par
  \nopagebreak
  \noindent\textit{Notas:} #4\par
}

% --- Separador de semana ------------------------------------
\newcommand{\semana}[2]{%
  \par\bigskip
  \noindent{\large\bfseries\color{tecnavy}Semana #1}\quad{\small\color{gris}#2}\par
  \vspace{-6pt}\noindent\rule{\textwidth}{0.4pt}\par
}

% --- Resumen final (se calcula solo) ------------------------
\newcommand{\resumen}{%
  \begin{center}
  \begin{tabular}{@{}lr@{}}
  \toprule
  Horas registradas       & \horasacum\,h \\
  Presupuesto total       & \presupuesto\,h \\
  Horas restantes         & \fpeval{\presupuesto - \horasacum}\,h \\
  Porcentaje consumido    & \fpeval{round(100*\horasacum/\presupuesto,1)}\,\% \\
  \bottomrule
  \end{tabular}
  \end{center}
}

\begin{document}

\begin{center}
{\LARGE\bfseries\color{tecnavy} AutoPipe --- Bitácora de avance}\\[4pt]
{\small Trabajo Final de Graduación --- Ricardo Antonio Borbón Mena}\\[2pt]
{\small Presupuesto: \presupuesto\,h (12 semanas $\times$ 28 h/semana)}
\end{center}

\vspace{6pt}

\section*{Bitácora}

% ============================================================
%  A partir de aquí: una linea \entrada{}{}{}{} por dia.
% ============================================================

\semana{1}{2026-08-03 a 2026-08-09 --- E1: análisis y diseño de la generalización}

\entrada{2026-08-04}{6}
{Verifiqué la línea base del repositorio migrado: 134 pruebas en verde, 97\% de
 cobertura y \texttt{ruff} sin hallazgos. Leí el anteproyecto completo y elaboré la
 hoja de ruta de 12 semanas (294\,h comprometidas, 42\,h de reserva). Creé esta
 bitácora y el pre-registro v0.1 del protocolo experimental.}
{Cerré cinco decisiones que la Sección~4.3 del anteproyecto no especifica: kNN como
 modelo base, parametrización fraccional del espacio de búsqueda, 3 semillas por
 condición, censura en 31 con chequeo de robustez, y Wilcoxon de una cola. Ninguna
 modifica los criterios de éxito. Trabajo en la rama
 \texttt{docs/planning-and-preregistration}; no se commitea directo en
 \texttt{develop}.}

\entrada{2026-08-05}{5}
{Escribí el ADR-001 con el diseño del contrato \texttt{fit/transform} y lo implementé
 en \texttt{BaseTransform} (\texttt{fit}, \texttt{transform}, \texttt{fit\_transform})
 y en \texttt{StatefulTransform}, con 30 pruebas nuevas. Suite en 164 pruebas, 97\% de
 cobertura, \texttt{ruff} limpio y ningún test preexistente modificado.}
{Hallazgo que cambió el diseño: el espacio de búsqueda comparte instancias de
 transformación entre candidatos, así que al haber estado hay que clonar el pipeline
 dentro de \texttt{evaluate\_pipeline}. Deuda detectada: \texttt{mypy} no corre
 (\texttt{python\_version} 3.10 contra stubs de numpy de 3.12). Pendientes de la
 Semana 1: los dos spikes (\texttt{enqueue\_trial} y \texttt{pymfe}) y el diseño de
 la interfaz de tarea.}

% ------------------------------------------------------------
%  Agregue las nuevas entradas debajo de esta línea.
% ------------------------------------------------------------

\newpage
\section*{Resumen de horas}

\resumen

\subsection*{Horas por semana: plan contra real}

Complete la columna \textit{Real} al cerrar cada semana. La columna \textit{Plan}
corresponde a las horas comprometidas de la hoja de ruta; la diferencia con las
28\,h semanales de capacidad es la reserva de contingencia.

\begin{center}
\begin{tabular}{@{}clrrl@{}}
\toprule
\textbf{Sem.} & \textbf{Etapas} & \textbf{Plan (h)} & \textbf{Real (h)} & \textbf{Observaciones} \\
\midrule
1  & E1                & 26 & --- & \\
2  & E1, E2            & 26 & --- & \\
3  & E2, E8            & 26 & --- & \\
4  & E2, E8            & 25 & --- & Hito: núcleo generalizado \\
5  & E3, E4, E8        & 25 & --- & \\
6  & E4, E8            & 25 & --- & \\
7  & E4, E5, E8        & 24 & --- & Hito: base y control de alcance \\
8  & E5, E8            & 24 & --- & \\
9  & E5, E6, E8        & 23 & --- & Hito: recomendador funcional \\
10 & E6, E7, E8        & 22 & --- & \\
11 & E6, E7, E8        & 22 & --- & \\
12 & E8                & 21 & --- & Hito: cierre del proyecto \\
\midrule
\multicolumn{2}{@{}l}{\textbf{Total comprometido}} & \textbf{289} & --- & \\
\multicolumn{2}{@{}l}{\textbf{Reserva de contingencia}} & \textbf{47} & --- & 14\,\% de \presupuesto\,h \\
\bottomrule
\end{tabular}
\end{center}

\subsection*{Decisiones abiertas}

Registro corto de decisiones que aún no están cerradas, para no perderlas entre
entradas diarias. Elimine el ítem cuando la decisión quede tomada (y anótela ese
día en la bitácora). Las decisiones ya cerradas sobre el protocolo experimental
viven en \texttt{docs/Preregistro.tex}, no aquí.

\begin{itemize}[leftmargin=1.4em]
  \item Subir de 3 a 5 semillas por condición: depende del costo de cómputo real
        medido en la semana 6.
  \item Tope de submuestreo para las meta-características del grupo de complejidad.
  \item Entrada o descarte del agrupamiento (\textit{clustering}) en el control de
        alcance de la semana 7.
\end{itemize}

\end{document}
